\section{Superincreasing Knapsack}
\label{sec:superincreasing}

Untuk memungkinkan proses dekripsi yang efisien, sistem Merkle--Hellman tidak
memakai barisan bobot sembarangan. Ia memakai barisan yang punya sifat khusus
sehingga persoalan knapsack yang tadinya sulit berubah menjadi pekerjaan yang
dapat diselesaikan sekali jalan, tanpa perlu mencoba-coba. Sifat khusus itu
disebut \emph{superincreasing}, dan barisan yang memilikinya adalah bahan baku
seluruh sistem.

\subsection{Barisan \textit{Superincreasing}}

Sebuah barisan bobot $\{w_1, w_2, \dots, w_n\}$ disebut \textbf{barisan
\textit{superincreasing}} jika setiap elemen di dalam barisan lebih besar
daripada jumlah semua elemen sebelumnya:

\begin{equation}
w_i > \sum_{j=1}^{i-1} w_j \quad \text{untuk semua } i = 2, \dots, n .
\label{eq:superincreasing}
\end{equation}

Syarat itu jauh lebih kuat daripada sekadar ``barisan menaik''. Sebagai
contoh, barisan $\{1, 2, 4, 9, 18, 36, 72, 144\}$ menaik sekaligus
superincreasing, sebab $9 > 1+2+4$, $18 > 1+2+4+9$, dan seterusnya. Sebaliknya
$\{1, 3, 4, 9, 15, 25\}$ menaik tetapi \emph{bukan} superincreasing, sebab
$4$ tidak lebih besar daripada $1 + 3 = 4$. Perbedaan itu tampak sepele, tetapi
justru dari situlah seluruh kemudahan dekripsi berasal.

Akibat langsung dari Persamaan~\eqref{eq:superincreasing} adalah bahwa elemen
terbesar selalu menguasai seluruh elemen di belakangnya. Jumlah semua elemen
selain $w_n$ paling banyak $w_n - 1$, sehingga tidak ada kombinasi elemen yang
lebih kecil yang dapat menandingi atau melampaui $w_n$. Kenyataan itu dipakai
oleh algoritma \textit{greedy} pada Subbagian~\ref{subsec:greedy-knapsack}.

\begin{table}[htbp]
    \centering
    \caption{Pemeriksaan dua barisan: satu superincreasing, satu bukan}
    \label{tab:barisan-superincreasing}
    \begin{tabularx}{\linewidth}{@{}c
                                   >{\raggedright\arraybackslash}X
                                   >{\raggedright\arraybackslash}X@{}}
    \toprule
    \textbf{Elemen} & \textbf{$\{1, 3, 6, 13, 27, 52\}$} &
    \textbf{$\{1, 3, 4, 9, 15, 25\}$} \\
    \midrule
    $w_1$ & $1$, elemen pertama & $1$, elemen pertama \\
    \addlinespace
    $w_2$ & $3 > 1$ \checkmark & $3 > 1$ \checkmark \\
    \addlinespace
    $w_3$ & $6 > 1 + 3 = 4$ \checkmark & $4 > 1 + 3 = 4$ \textbf{gagal} \\
    \addlinespace
    $w_4$ & $13 > 1 + 3 + 6 = 10$ \checkmark & tidak perlu diperiksa lagi \\
    \addlinespace
    $w_5$ & $27 > 1 + 3 + 6 + 13 = 23$ \checkmark & --- \\
    \addlinespace
    $w_6$ & $52 > 1 + 3 + 6 + 13 + 27 = 50$ \checkmark & --- \\
    \addlinespace
    Kesimpulan & \textbf{superincreasing} & \textbf{bukan} \\
    \bottomrule
    \end{tabularx}
\end{table}

Satu kekeliruan yang sering terjadi adalah memeriksa hanya elemen terakhirnya.
Tabel~\ref{tab:barisan-superincreasing} memperlihatkan bahwa syarat itu harus
dipenuhi oleh \emph{setiap} elemen, dan pada barisan kedua pelanggaran terjadi
sudah pada elemen ketiga. Pada barisan $\{1, 3, 4, 9, 15, 25\}$ misalnya,
$25 > 1+3+4+9+15 = 32$ juga gagal --- jadi barisan itu melanggar syarat dua
kali, bukan sekali.

\subsection{Mengapa Algoritma \textit{Greedy} Selalu Berhasil}
\label{subsec:greedy-knapsack}

Kemudahan dekripsi superincreasing knapsack bertumpu pada satu pengamatan
sederhana, yang dapat dinyatakan sebagai lema berikut.

\begin{lemma}[Keunikan pilihan elemen terbesar]
\label{lem:greedy-superincreasing}
Misalkan $W = \{w_1, \dots, w_n\}$ barisan superincreasing dan $M$ adalah
sebuah bilangan dengan $0 \le M \le \sum_{i=1}^{n} w_i$. Bila
$M = \sum_{i=1}^{n} b_i w_i$ untuk suatu vektor biner, maka $b_n = 1$ jika dan
hanya jika $M \ge w_n$.
\end{lemma}

Alasannya sebagai berikut. Andaikan $b_n = 0$. Maka seluruh bobot yang
terpilih berasal dari $\{w_1, \dots, w_{n-1}\}$, sehingga
$M \le \sum_{j=1}^{n-1} w_j < w_n$ menurut
Persamaan~\eqref{eq:superincreasing}. Jadi $M < w_n$, yang berarti
$M \ge w_n$ mustahil terjadi bila $b_n = 0$. Sebaliknya, bila $b_n = 1$, maka
$M \ge w_n$ jelas terpenuhi karena seluruh $w_i$ positif.

Lema~\ref{lem:greedy-superincreasing} itulah yang membuat dekripsi menjadi
pekerjaan sekali jalan. Setelah $b_n$ ditentukan, sisa bobot
$M - b_n w_n$ harus dibentuk oleh $\{w_1, \dots, w_{n-1}\}$ yang juga
superincreasing, sehingga lema yang sama dapat diterapkan lagi. Dengan begitu
seluruh bit dapat ditentukan dari belakang ke depan tanpa pernah menebak dan
tanpa pernah kembali memperbaiki keputusan yang sudah diambil --- itulah yang
dimaksud dengan algoritma \emph{greedy}.

\subsection{Prosedur Dekripsi Mudah}

Berbeda dengan \textit{knapsack} umum, \textit{superincreasing knapsack} dapat
dipecahkan dalam waktu linier $O(n)$ dengan algoritma \textit{greedy}:

\begin{enumerate}
    \item Bandingkan bobot total $M$ dengan bobot terbesar $w_n$.
    \item Jika $w_n \le M$, maka $b_n = 1$ dan kurangi $M$ dengan $w_n$ ($M = M - w_n$). Jika tidak, $b_n = 0$.
    \item Ulangi proses tersebut untuk bobot $w_{n-1}, w_{n-2}, \dots, w_1$.
    \item Jika pada akhir proses $M = 0$, maka solusi ditemukan.
\end{enumerate}

Perhatikan langkah 4. Ia bukan sekadar pengaman, melainkan penentu: bila setelah
seluruh bobot diperiksa sisanya belum nol, maka $M$ tadi memang tidak dapat
dibentuk oleh barisan itu, dan jawaban yang diperoleh bukan solusi yang sah.
Pada pemakaian kriptografis hal ini tidak boleh terjadi, sebab kriptogram selalu
dihasilkan dari sebuah blok pesan yang sah.

Contoh lengkap penerapan keempat langkah itu disajikan pada
Contoh~\ref{ex:greedy-superincreasing}.

\subsection{Mengapa Kemudahan Itu Menjadi Kelemahan}

Sifat yang membuat dekripsi mudah justru membuat barisan superincreasing tidak
layak dipakai langsung sebagai kunci. Persoalannya bukan pada kecepatannya,
melainkan pada kenyataan bahwa siapa pun --- termasuk penyadap --- dapat
menjalankan algoritma yang sama. Bila barisan bobot dirahasiakan sekalipun,
penyerang tidak perlu mengetahuinya: ia cukup menebak, dan tebakannya dapat
diperiksa seketika, sebab nilai $M$ tersedia di tangannya.

Karena itu Merkle dan Hellman menempuh jalan lain. Barisan superincreasing
dipertahankan sebagai kunci privat --- tetap mudah untuk pemiliknya --- tetapi
tidak pernah dikirimkan atau disiarkan. Yang disiarkan adalah wujud lain dari
barisan yang sama, yang telah diubah sedemikian rupa sehingga sifat
superincreasing-nya tidak lagi terlihat. Cara mengubahnya, dan mengapa
perubahan itu tidak benar-benar menyembunyikan apa pun, dibahas pada bagian
berikutnya.
